\documentclass[a4paper,12pt]{article}
\usepackage{ctex}
\usepackage{geometry}
\usepackage{enumitem}
\usepackage{amsmath}

\geometry{left=2.5cm, right=2.5cm, top=2.5cm, bottom=2.5cm}

\title{\textbf{第四章\ 函数}}
\author{}
\date{}

\begin{document}

\maketitle

\section*{一、选择题}

\begin{enumerate}
	\item 以下叙述中正确的是
	      \begin{enumerate}
		      \item[A)] 构成C程序的基本单位是函数
		      \item[B)] 可以在一个函数中定义另一个函数
		      \item[C)] main（）函数必须放在其他函数之前
		      \item[D)] C函数定义的格式是 K\&R格式
	      \end{enumerate}
	      \vspace{0.3cm}

	\item 一个C语言程序是由
	      \begin{enumerate}
		      \item[A)] 一个主程序和若干子程序组成
		      \item[B)] 函数组成
		      \item[C)] 若干过程组成
		      \item[D)] 若干子程序组成
	      \end{enumerate}
	      \vspace{0.3cm}

	\item C语言程序的基本单位是
	      \begin{enumerate}
		      \item[A)] 程序行
		      \item[B)] 语句
		      \item[C)] 函数
		      \item[D)] 字符
	      \end{enumerate}
	      \vspace{0.3cm}

	\item 在C语言中，以下哪个选项表示函数返回类型为 \texttt{void}？
	      \begin{enumerate}
		      \item[A)] \texttt{int}
		      \item[B)] \texttt{float}
		      \item[C)] \texttt{char}
		      \item[D)] \texttt{void}
	      \end{enumerate}
	      \vspace{0.3cm}

	\item 在C语言中，以下哪个选项表示函数参数列表为空？
	      \begin{enumerate}
		      \item[A)] \texttt{int}
		      \item[B)] \texttt{void}
		      \item[C)] \texttt{char}
		      \item[D)] \texttt{()}
	      \end{enumerate}
	      \vspace{0.3cm}

	\item 在C语言中，以下哪个选项表示函数的返回值类型为整型？
	      \begin{enumerate}
		      \item[A)] \texttt{int}
		      \item[B)] \texttt{float}
		      \item[C)] \texttt{char}
		      \item[D)] \texttt{void}
	      \end{enumerate}
	      \vspace{0.3cm}

	\item 在C语言中，以下哪个选项表示函数的返回值类型为浮点型？
	      \begin{enumerate}
		      \item[A)] \texttt{int}
		      \item[B)] \texttt{float}
		      \item[C)] \texttt{char}
		      \item[D)] \texttt{void}
	      \end{enumerate}
	      \vspace{0.3cm}

	\item 在C语言中，以下哪个选项表示函数的返回值类型为字符型？
	      \begin{enumerate}
		      \item[A)] \texttt{int}
		      \item[B)] \texttt{float}
		      \item[C)] \texttt{char}
		      \item[D)] \texttt{void}
	      \end{enumerate}
	      \vspace{0.3cm}

	\item 下列对C程序结构的说法不正确的是
	      \begin{enumerate}
		      \item[A)] C程序是由一系列函数构成的
		      \item[B)] C程序可以有多个main（）函数
		      \item[C)] C程序中函数名不可以和变量名相同
		      \item[D)] C程序中可以定义函数
	      \end{enumerate}
	      \vspace{0.3cm}

	\item C语言规定，程序中各函数之间
	      \begin{enumerate}
		      \item[A)] 既允许直接递归调用也允许间接递归调用
		      \item[B)] 不允许直接递归调用也不允许间接递归调用
		      \item[C)] 允许直接递归调用不允许间接递归调用
		      \item[D)] 不允许直接递归调用允许间接递归调用
	      \end{enumerate}
	      \vspace{0.3cm}

	\item 设有如下的函数 \texttt{exam(float x)\{ printf("\%f",x*x); \}}，则函数的类型为\_\_\_\_\_\_。
	      \begin{enumerate}
		      \item[A)] 与参数x的类型相同
		      \item[B)] 是void
		      \item[C)] 是int
		      \item[D)] 无法确定
	      \end{enumerate}
	      \vspace{0.3cm}

	\item 以下关于宏的叙述中正确的是
	      \begin{enumerate}
		      \item[A)] 宏名必须用大写字母表示
		      \item[B)] 宏定义必须位于源程序中所有语句之前
		      \item[C)] 宏替换没有数据类型限制
		      \item[D)] 宏调用比函数调用耗费时间
	      \end{enumerate}
	      \vspace{0.3cm}

	\item 宏替换\_\_\_\_\_\_占用运行时间。
	      \begin{enumerate}
		      \item[A)] 不
		      \item[B)] 轻微
		      \item[C)] 大量
		      \item[D)] 视情况而定
	      \end{enumerate}
	      \vspace{0.3cm}

	\item 以下叙述中正确的是
	      \begin{enumerate}
		      \item[A)] C程序的基本组成单位是语句
		      \item[B)] C程序中的每一行只能写一条语句
		      \item[C)] 简单C语句必须以分号结束
		      \item[D)] C语句必须在一行内写完
	      \end{enumerate}
	      \vspace{0.3cm}

	\item 若a是int型变量，且a=5，则表达式 \texttt{(a+100)\%2+a/2} 的值为\_\_\_\_\_\_。
	      \begin{enumerate}
		      \item[A)] 52
		      \item[B)] 51
		      \item[C)] 53
		      \item[D] 50
	      \end{enumerate}
	      \vspace{0.3cm}

	\item 下面函数是求两个整型参数a和b的最小公倍数，空白处应填\_\_\_\_\_\_。
	      \begin{verbatim}
int f2(int a, int b)
{
    int i=2, p=1;
    do {
        while(a%i==0 && ____________) {
            p*=i; a/=i; b/=i;
        }
        ____________;
    }while(a>=i && ____________);
    return p*a*b;
}
    \end{verbatim}
	      \begin{enumerate}
		      \item[A)] b%i==0 / i++ / b>=i
		      \item[B)] b\%i==0 / ++i / b>i
		      \item[C)] b\%i / i++ / b>=i
		      \item[D] b\%i==0 / i++ / b>i
	      \end{enumerate}
	      \vspace{0.3cm}

	\item 主函数调用一个fun函数将字符串逆序，空白处应填\_\_\_\_\_\_。
	      \begin{verbatim}
#include <stdio.h>
#include <string.h>
____________;
int main() {
    char s[80];
    scanf("%s", s);
    ____________;
    printf("逆序后的字符串:%s\n", s);
    return 0;
}
void fun(char ss[]) {
    int n = strlen(ss);
    for(int i = 0; ____________; i++) {
        char c = ss[i];
        ss[i] = ss[n-1-i];
        ss[n-1-i] = c;
    }
}
    \end{verbatim}
	      \begin{enumerate}
		      \item[A)] void fun(char ss[]) / fun(s) / i<n/2
		      \item[B)] char fun(char ss[]) / fun(s) / i<n
		      \item[C)] void fun(char ss[]) / s=fun(s) / i<n/2
		      \item[D] void fun() / fun(s) / i<n/2
	      \end{enumerate}
	      \vspace{0.3cm}

	\item 有以下程序
	      \begin{verbatim}
#include <stdio.h>
void fun(char *c, int d)
{
    *c=*c+1; d=d+1;
    printf("%c, %c, ", *c, d);
}
int main()
{
    char a='A', b='a';
    fun(&b, a);
    printf("%c, %c\n", a, b);
    return 0;
}
    \end{verbatim}
	      程序运行后的输出结果是\_\_\_\_\_\_。
	      \begin{enumerate}
		      \item[A)] B, b, A, a
		      \item[B)] a, B, B, b
		      \item[C)] A, b, B, a
		      \item[D)] b, B, a, B
	      \end{enumerate}
	      \vspace{0.3cm}

	\item
	      \begin{verbatim}
#include <stdio.h>
int fun(int a,int b,int*cn,int*dn){
    *cn=a*b+b*b;
    *dn=a*a-b*b;
    a=5; b=6;
}
int main(){
    int a=2,b=3,c=5,d=6;
    fun(a,b,&c,&d);
    printf("a=%d,b=%d,c=%d\n",a,b,c,d);
    return 0;
}
    \end{verbatim}
	      输出\_\_\_\_\_\_。
	      \begin{enumerate}
		      \item[A)] a=2,b=3,c=20,d=-5
		      \item[B)] a=2,b=3,c=15,d=5
		      \item[C] a=5,b=6,c=20,d=-5
		      \item[D] a=5,b=6,c=15,d=5
	      \end{enumerate}
	      \vspace{0.3cm}

	\item
	      \begin{verbatim}
#include <stdio.h>
int fun(int x){
    static y=2;
    y++;
    x+=y;
    return x;
}
int main(){
    int k;
    k=fun(3);
    printf("%d,%d\n",k,fun(k));
    return 0;
}
    \end{verbatim}
	      输出\_\_\_\_\_\_。
	      \begin{enumerate}
		      \item[A)] 6,9
		      \item[B)] 6,8
		      \item[C] 5,7
		      \item[D] 5,8
	      \end{enumerate}
	      \vspace{0.3cm}

	\item
	      \begin{verbatim}
int f(int n){
    if(n==0||n==1) return 1;
    else return f(n-2)+2*f(n-1);
}
int main(){
    int n=5;
    printf("%d",f(n));
    return 0;
}
    \end{verbatim}
	      输出\_\_\_\_\_\_。
	      \begin{enumerate}
		      \item[A)] 11
		      \item[B] 13
		      \item[C] 15
		      \item[D] 17
	      \end{enumerate}
	      \vspace{0.3cm}

	\item
	      \begin{verbatim}
main()
{ void fun();
  float x,y;
  x=10;
  fun(x,&y);
  printf("result:%.0f,%.0f\n",x,y);
}
void fun(x,y)
float x,*y;
{ *y=x*x;
}
    \end{verbatim}
	      输出\_\_\_\_\_\_。
	      \begin{enumerate}
		      \item[A)] result:10,100
		      \item[B)] result:10,10
		      \item[C] result:100,100
		      \item[D] result:100,10
	      \end{enumerate}
	      \vspace{0.3cm}

	\item 下面程序的功能是\_\_\_\_\_\_。
	      \begin{verbatim}
void ss(char *s1,char*s2){
    while(*s1!='\0') s1++;
    while(*s2!='\0'){
        *s1=*s2;
        s1++; s2++;
    }
    *s1='\0';
}
    \end{verbatim}
	      \begin{enumerate}
		      \item[A)] 字符串连接
		      \item[B)] 字符串比较
		      \item[C] 字符串复制
		      \item[D] 字符串查找
	      \end{enumerate}
	      \vspace{0.3cm}

	\item C语言规定，函数返回值的类型是由\_\_\_\_\_\_决定的。
	      \begin{enumerate}
		      \item[A)] return语句中的表达式类型
		      \item[B)] 函数定义时的返回类型
		      \item[C] 调用时的实参类型
		      \item[D] 以上都不对
	      \end{enumerate}
	      \vspace{0.3cm}

	\item 如果在一个函数中的复合语句中定义了一个变量，则该变量\_\_\_\_\_\_。
	      \begin{enumerate}
		      \item[A)] 只在该复合语句中有效
		      \item[B)] 在该函数中有效
		      \item[C] 在整个程序中有效
		      \item[D] 在程序运行期间一直有效
	      \end{enumerate}
	      \vspace{0.3cm}

	\item C语言中的函数\_\_\_\_\_\_递归定义。
	      \begin{enumerate}
		      \item[A)] 可以
		      \item[B)] 不可以
		      \item[C] 只能直接
		      \item[D] 只能间接
	      \end{enumerate}
	      \vspace{0.3cm}

	\item main函数\_\_\_\_\_\_参数。
	      \begin{enumerate}
		      \item[A)] 可以有
		      \item[B)] 不可以有
		      \item[C] 只能有一个
		      \item[D] 只能有两个
	      \end{enumerate}
	      \vspace{0.3cm}
\end{enumerate}

\section*{二、填空题}

\begin{enumerate}
	\item 在C语言中，一个函数一般由两个部分组成，它们分别是\_\_\_\_\_\_和\_\_\_\_\_\_。
	      \vspace{0.3cm}

	\item 下面函数是求两个整型参数a和b的最小公倍数，请在空格处填上正确的内容。
	      \begin{verbatim}
int f2(int a, int b)
{
    int i=2, p=1;
    do {
        while(a%i==0 && ____________) {
            p*=i; a/=i; b/=i;
        }
        ____________;
    }while(a>=i && ____________);
    return p*a*b;
}
    \end{verbatim}
	      \vspace{0.3cm}

	\item 主函数调用一个fun函数将字符串逆序，请在空格处填上正确的内容。
	      \begin{verbatim}
#include <stdio.h>
#include <string.h>
____________;
int main() {
    char s[80];
    scanf("%s", s);
    ____________;
    printf("逆序后的字符串:%s\n", s);
    return 0;
}
void fun(char ss[]) {
    int n = strlen(ss);
    for(int i = 0; ____________; i++) {
        char c = ss[i];
        ss[i] = ss[n-1-i];
        ss[n-1-i] = c;
    }
}
    \end{verbatim}
	      \vspace{0.3cm}

	\item 本函数用对分查找法，在已按字母次序从小到大排序的字符数组list中查找字符c，若c在数组中，函数返回字符c在数组中的下标，否则返回-1。请填空：
	      \begin{verbatim}
int search(char list[],char c,int len){
    int low,high,k;
    low=0; high=len-1;
    while( ____________ ) {
        k=(low+high)/2;
        if( ____________ ) return k;
        else if( ____________ ) high=k-1;
        else low=k+1;
    }
    return -1;
}
    \end{verbatim}
	      \vspace{0.3cm}

	\item 函数mycmp(char *s, char *t)的功能是比较字符串s和t的大小，当s等于t时返回0，否则返回s和t的第一个不同字符的ASCII码的差值。请填空：
	      \begin{verbatim}
int mycmp(char*s,char*t){
    while(*s==*t){
        if( ____________ ) return 0;
        ++s; ++t;
    }
    return( ____________ );
}
    \end{verbatim}
	      \vspace{0.3cm}

	\item 下面函数的功能是\_\_\_\_\_\_。
	      \begin{verbatim}
int max(int x, int y) {
    return (x > y) ? x : y;
}
    \end{verbatim}
	      \vspace{0.3cm}

	\item 下面函数的功能是\_\_\_\_\_\_。
	      \begin{verbatim}
int isPrime(int n) {
    if (n <= 1) return 0;
    for (int i = 2; i * i <= n; i++)
        if (n % i == 0) return 0;
    return 1;
}
    \end{verbatim}
	      \vspace{0.3cm}

	\item 下面函数的功能是\_\_\_\_\_\_。
	      \begin{verbatim}
int factorial(int n) {
    if (n <= 1) return 1;
    return n * factorial(n - 1);
}
    \end{verbatim}
	      \vspace{0.3cm}

	\item 下面函数的功能是\_\_\_\_\_\_。
	      \begin{verbatim}
int fibonacci(int n) {
    if (n <= 2) return 1;
    return fibonacci(n - 1) + fibonacci(n - 2);
}
    \end{verbatim}
	      \vspace{0.3cm}
\end{enumerate}

\section*{三、阅读程序题}

\begin{enumerate}
	\item
	      \begin{verbatim}
#include <stdio.h>
void fun(char *c, int d)
{
    *c=*c+1; d=d+1;
    printf("%c, %c, ", *c, d);
}
int main()
{
    char a='A', b='a';
    fun(&b, a);
    printf("%c, %c\n", a, b);
    return 0;
}
    \end{verbatim}
	      程序运行后的输出结果是\_\_\_\_\_\_。
	      \vspace{0.3cm}

	\item
	      \begin{verbatim}
#include <stdio.h>
int fun(int a,int b,int*cn,int*dn){
    *cn=a*b+b*b;
    *dn=a*a-b*b;
    a=5; b=6;
}
int main(){
    int a=2,b=3,c=5,d=6;
    fun(a,b,&c,&d);
    printf("a=%d,b=%d,c=%d\n",a,b,c,d);
    return 0;
}
    \end{verbatim}
	      输出\_\_\_\_\_\_。
	      \vspace{0.3cm}

	\item
	      \begin{verbatim}
#include <stdio.h>
int fun(int x){
    static y=2;
    y++;
    x+=y;
    return x;
}
int main(){
    int k;
    k=fun(3);
    printf("%d,%d\n",k,fun(k));
    return 0;
}
    \end{verbatim}
	      输出\_\_\_\_\_\_。
	      \vspace{0.3cm}

	\item
	      \begin{verbatim}
int f(int n){
    if(n==0||n==1) return 1;
    else return f(n-2)+2*f(n-1);
}
int main(){
    int n=5;
    printf("%d",f(n));
    return 0;
}
    \end{verbatim}
	      输出\_\_\_\_\_\_。
	      \vspace{0.3cm}

	\item
	      \begin{verbatim}
main()
{ void fun();
  float x,y;
  x=10;
  fun(x,&y);
  printf("result:%.0f,%.0f\n",x,y);
}
void fun(x,y)
float x,*y;
{ *y=x*x;
}
    \end{verbatim}
	      输出\_\_\_\_\_\_。
	      \vspace{0.3cm}

	\item 写出下面程序的功能：
	      \begin{verbatim}
void ss(char *s1,char*s2){
    while(*s1!='\0') s1++;
    while(*s2!='\0'){
        *s1=*s2;
        s1++; s2++;
    }
    *s1='\0';
}
    \end{verbatim}
	      \vspace{0.3cm}
\end{enumerate}

\section*{四、编程题}

\begin{enumerate}
	\item 请将下列一组数据读入到S数组中，并从中找出最小的值并输出。\\
	      30, 56, 88, 45, 100, 20
	      \vspace{2cm}

	\item 请将下列给出的字符串读入到ss数组中，并输出该字符串。\\
	      Student and Teacher
	      \vspace{2cm}

	\item 编写一个程序，从键盘输入一个整数，判断该整数是否为素数。
	      \vspace{2cm}

	\item 编写一个程序，计算1到n的阶乘之和。
	      \vspace{2cm}

	\item 编写一个程序，输入一个字符串，将其转换为大写字母并输出。
	      \vspace{2cm}

	\item 编写一个程序，输入一个整数，将其转换为二进制字符串并输出。
	      \vspace{2cm}

	\item 从键盘上任意输入一个字符串，统计字符串中大小写英文字母出现的次数。
	      \vspace{2cm}

	\item 编写程序求某三位数的平方和（8分）。
	      \vspace{2cm}

	\item 编写程序求数列1,1,2,3,5,8,13,...的前100项的和及平均值（12分）。
	      \vspace{2cm}
\end{enumerate}

\section*{五、改错题}

\begin{enumerate}
	\item 找出下面程序的错误并改正：
	      \begin{verbatim}
#include<stdio.h>
char a="Beijing";
int main()
{
    printf("%s is one city in china.\n",a);
    p1();
    p2();
    return 0;
}
int p1()
{
    printf("%s is one of the biggest city",a);
    return;
}
int p2()
{
    printf("in the world.\n");
    return;
}
    \end{verbatim}
	      错误的行是：\_\_\_\_\_\_ 改为：\_\_\_\_\_\_
	      \vspace{0.3cm}

	\item 求 $\sum_{k=1}^{100} k + \sum_{k=1}^{50} k^2 + \sum_{k=1}^{10} k$，找出下面程序的错误并改正：
	      \begin{verbatim}
#include <stdio.h>
int main()
{
    int n1=100,n2=50,n3=10;
    int k;
    float s1=0,s2=0,s3=0;
    for(k=1;k<=n1;k++);
        s1=s1+k;
    for(k=1;k<=n2;k++);
        s2=s2+k*k;
    for(k=1;k<=n3;k++);
        s3=s3+k/10;
    printf("total=%8.2f\n",s1+s2+s3);
    return 0;
}
    \end{verbatim}
	      错误的行是：\_\_\_\_\_\_ 改为：\_\_\_\_\_\_
	      \vspace{0.3cm}

	\item 将给定的字符串写入指定的文件中去，找出下面程序的错误并改正：
	      \begin{verbatim}
#include <stdio.h>
int main()
{
    int i;
    FILE*fp;
    char a[][8]={"Turbo C", "BASIC", "FORTRAN", "COBOL"};
    if((fp=fopen("myfile.txt", "w"))==NULL)
    {
        printf("can't creat the file.\n");
        return 1;
    }
    for(i=0;i<4;i++)
    {
        fputs(a[i],fp);
        fputs("\n",fp);
    }
    fclose(fp);
    return 0;
}
    \end{verbatim}
	      错误的行是：\_\_\_\_\_\_ 改为：\_\_\_\_\_\_
	      \vspace{0.3cm}
\end{enumerate}

\end{document}
